% =========================================================
% 5.2  Error-Conditioned Pedagogical Signals in the Constructed Preference Dataset
% ---------------------------------------------------------
% Draft for chapters/05_results.tex (Research Question 1). Replaces the section skeleton.
% Figures: result_5_2/fig_5_6_associations, fig_5_7_context_ranking, fig_5_8_posterior_shift,
%          fig_5_11_margin_by_pair, fig_5_12_model_composition, and the appendix figure fig_A_contexts (appendix_A_contexts.tex);
%          all built by result_5_2.ipynb at the repository root, to be copied into
%          __Thesis_writing/TUM_Template/figures/ by hand.
% Every number below is read off one of the figures, the table, or the appendix figure.
% Section 4.5.3 (DPO subsets) is not referenced; the margin paragraph describes the pairs only.
% =========================================================
\section{Error-Conditioned Pedagogical Signals in the Constructed Preference Dataset}\label{sec:preference_signal_results}
This section addresses Research Question~1. Section~\ref{subsec:observed_associations} reports the associations between student error types and first-turn tutor moves that the scoring method of Section~\ref{subsec:error_conditioned_preference_scoring} builds on. Section~\ref{subsec:associations_in_pairs} then examines whether these associations are reflected in the 24,298 preference pairs that Section~\ref{subsec:preference_pair_construction} constructs from the 48,839 candidate responses of the 693 training contexts. All figures are computed from the saved candidate table and the saved preference pairs; no response was generated or re-scored for this section.

% ---------------------------------------------------------
\subsection{Observed Error--Pedagogy Associations in the Source Dataset}\label{subsec:observed_associations}
Figure~\ref{fig:associations} shows the first-turn co-occurrence table of Table~\ref{tab:first_turn_error_move_cooccurrence} in two views: on the left, the observed counts, that is, the number of training contexts whose first tutor turn carries each move, per error type; on the right, the PPMI utilities of Table~\ref{tab:first_turn_error_move_ppmi} that the scoring method derives from these counts.

    \begin{figure}[!t]
        \centering
        \includegraphics[width=\textwidth]{figures/fig_5_6_associations.pdf}
        \caption{Observed associations between student error types and first-turn tutor moves in 693 dialogue contexts from the training split of the aligned dataset. Left: number of contexts whose first tutor turn carries each move, per error type (the counts of Table~\ref{tab:first_turn_error_move_cooccurrence}), with the number of contexts per error type in brackets and the number of first turns per move below each move name. Right: the corresponding utility \(u(m,e)=\operatorname{PPMI}(m,e)\) of Table~\ref{tab:first_turn_error_move_ppmi}. In both panels, darker cells hold larger values.}
        \label{fig:associations}
    \end{figure}

    \paragraph{Co-occurrence and PPMI Patterns}
    \textit{Generic} is the most frequent first move for every error type: it accounts for 529 of the 693 first tutor turns, from 27 of the 34 turns of \textit{unit conversion error} to 160 of the 199 turns of \textit{misunderstanding of the question}. \textit{Probing} follows with 108 first turns, \textit{Focus} with 47, and \textit{Telling} with 9. The right panel shows how PPMI turns this frequency picture into a different one. Of the 28 cells, 14 have a positive utility and 14 are zero, and every error type has between one and three moves with a positive utility: one for \textit{numerical calculation error} (\textit{Focus}) and for \textit{misunderstanding of the question} (\textit{Generic}), two for \textit{missing or extra quantity}, \textit{none of the above}, and \textit{reached the correct solution but proceeded further}, and three for \textit{missing or incorrect factual knowledge} and \textit{unit conversion error}. Because the score \(s(r,e)\) is the posterior-weighted sum of these utilities, the zero cells define what the scoring does not reward: posterior mass on a zero-utility move adds nothing to the score, and, since the posterior sums to one, it is mass that the moves with a positive utility do not receive. \textit{Generic} receives a positive utility only for the four error types in which it occurs more often than its overall rate of 529 in 693 (\textit{missing or extra quantity}, \textit{misunderstanding of the question}, \textit{none of the above}, and \textit{unit conversion error}), and these are the smallest utilities in the table (0.04--0.07). The move with the largest utility differs between error types: \textit{Focus} for \textit{numerical calculation error} (0.90), \textit{none of the above} (0.30), and \textit{reached the correct solution but proceeded further} (0.27); \textit{Telling} for \textit{missing or extra quantity} (0.47), \textit{missing or incorrect factual knowledge} (1.67), and \textit{unit conversion error} (1.18); and \textit{Generic} for \textit{misunderstanding of the question} (0.07). \textit{Probing} has a positive utility for three error types (0.18--0.47) but is never the largest.

    \paragraph{Strength and Support of the Associations}
    The left panel also shows the observations behind each utility. The three largest utilities, all for \textit{Telling}, rest on three, four, and one first tutor turns out of 167, 97, and 34 contexts, respectively. The utility of \textit{Focus} for \textit{numerical calculation error} (0.90) rests on 11 of 87 contexts, and the three positive \textit{Probing} utilities on 21, 9, and 6 first turns. The utilities are therefore ordered differently from the counts: the largest values in the right panel come from the rarest move in the left panel. This property of the utilities is taken up again in Chapter~\ref{chapter:Discussions}.

% ---------------------------------------------------------
\subsection{Examining Error--Pedagogy Associations in Chosen Responses}\label{subsec:associations_in_pairs}
The scoring of Section~\ref{subsec:error_conditioned_preference_scoring} gives every candidate of a context a score, and the pairing rule of Section~\ref{subsec:preference_pair_construction} splits the ranked candidates of the context into a chosen half and a rejected half. This subsection follows the construction at four levels: two contexts in detail (Figure~\ref{fig:context_ranking} and Table~\ref{tab:example_pairs}), one context for each of the seven error types (Figure~\ref{fig:app_contexts} in Appendix~\ref{app:contexts}), the mean posterior of the chosen and of the rejected responses over all 24,298 pairs (Figure~\ref{fig:posterior_shift} in Appendix~\ref{app:posterior_comparison}), and the score margins of all pairs (Figure~\ref{fig:margin_by_pair}). In each of them, the zero cells of Figure~\ref{fig:associations} are as visible as the positive ones.

    \paragraph{Scoring of a Preference Pair}
    Figure~\ref{fig:context_ranking} follows two contexts through the scoring and pairing of Sections~\ref{subsec:error_conditioned_preference_scoring} and~\ref{subsec:preference_pair_construction}: one with a \textit{numerical calculation error} (context 609, 69 candidates, positive utility only for \textit{Focus}) and one with \textit{missing or incorrect factual knowledge} (context 150, 72 candidates, utilities of 0.28, 0.47, and 1.67 for \textit{Focus}, \textit{Probing}, and \textit{Telling}). The lower panels show the posterior of every candidate, the upper panels the resulting score, the ranking, and the first three pairs. For the \textit{numerical calculation error}, the three other utilities are zero, so the score is 0.90 times the \textit{Focus} posterior and every unit of posterior mass on \textit{Probing} or \textit{Telling} lowers it: the 34 chosen candidates all have a \textit{Focus} posterior above 0.87 and scores between 0.79 and 0.84, the score falls below 0.7 only from rank 50 onwards, and the last ten ranks are the ten candidates whose posterior is dominated by \textit{Telling} (0.50--0.70), with scores of 0.21--0.40. For \textit{missing or incorrect factual knowledge} the order is reversed, because \textit{Telling} now carries the largest utility and only \textit{Generic} is zero. The twelve candidates with a dominant \textit{Telling} posterior (0.77--0.88) occupy the first twelve ranks with scores of 1.35--1.50, the candidates that mix \textit{Focus} with some \textit{Telling} mass follow, and the 36 rejected candidates are \textit{Focus} responses with a \textit{Focus} posterior of 0.76--0.93 and scores of 0.31--0.37.

    Table~\ref{tab:example_pairs} lists the second pair (\(j=2\)) of both contexts. In both pairs, one response asks the student which quantity to recompute and the other gives the complete calculation. The scoring puts the question on the chosen side and the complete calculation on the rejected side for the \textit{numerical calculation error} (scores 0.83 against 0.24), and the complete calculation on the chosen side and the question on the rejected side for \textit{missing or incorrect factual knowledge} (1.48 against 0.31). The same two kinds of response thus receive opposite preferences under the two error types, which is the intended behavior of the utility \(u(m,e)\).

    \begin{figure}[!t]
        \centering
        \includegraphics[width=\textwidth]{figures/fig_5_7_context_ranking.pdf}
        \caption{All candidates of two training contexts, ranked by the error-conditioned score \(s(r,e)\). Left: \textit{numerical calculation error} (context 609, \(n=69\)); right: \textit{missing or incorrect factual knowledge} (context 150, \(n=72\)). The utility vector under each title lists \(u(m,e)\) in the order \textit{Focus}, \textit{Generic}, \textit{Probing}, \textit{Telling}. Top: the score of every candidate; the blue half is chosen, the orange half rejected, and the grey middle candidate is omitted when \(n\) is odd. The arcs mark the first three pairs, which join rank \(j\) with rank \(n-j+1\). Bottom: the classifier posterior \(p_{\theta}(m\mid r)\) of the same candidates, stacked per move.}
        \label{fig:context_ranking}
    \end{figure}

    % --- Table: generated by result_5_2.ipynb -> result_5_2/text/table_example_pairs.tex ---
    \begin{table}[!t]
        \centering
        \footnotesize
        \setlength{\tabcolsep}{3pt}
        \renewcommand{\arraystretch}{1.15}
        \captionsetup{skip=4pt}
        \begin{tabularx}{\linewidth}{@{}l>{\raggedright\arraybackslash}Xlr@{}}
            \toprule
            Side & Tutor response (generation model) & \mbox{$p_{\theta}(m \mid r)$} & \mbox{$s(r, e)$} \\
            \midrule
            \multicolumn{4}{@{}l}{\textit{Numerical calculation error} (context 609), $u(m, e)$ = (0.90, 0, 0, 0)} \\
            Chosen & let's look closely at the bridge trolls: the problem says "6 less than four times that number," which means we need to multiply first before subtracting 6. can you try calculating the bridge trolls again using that rule? (gemini-3.5-flash-lite) & \mbox{(0.93, 0.00, 0.05, 0.02)} & 0.83 \\
            Rejected & “6 less than four times 6” means \(4 \times 6 - 6 = 18\) trolls under the bridge, not \(4 \times 6 + 6\). half of 18 is 9, so the total is \(6 + 18 + 9 = 33\) trolls. (gpt-5.6-terra) & \mbox{(0.26, 0.00, 0.07, 0.67)} & 0.24 \\
            \midrule
            \multicolumn{4}{@{}l}{\textit{Missing or incorrect factual knowledge} (context 150), $u(m, e)$ = (0.28, 0, 0.47, 1.67)} \\
            Chosen & the half-off offer must be used on the cheaper original pair, so the \$40 shoes cost \$20 and the \$60 shoes stay \$60, for a subtotal of \$80. a fourth of \$80 is \$20, so james pays \$80 - \$20 = \textbf{\$60}. (gpt-5.6-terra) & \mbox{(0.09, 0.01, 0.03, 0.86)} & 1.48 \\
            Rejected & you correctly found the prices of both pairs of shoes, but let's re-read carefully which pair is supposed to be "half off." which of the two prices, \$40 or \$60, is the cheaper one that gets the discount? (gemini-3.5-flash-lite) & \mbox{(0.88, 0.00, 0.11, 0.01)} & 0.31 \\
            \bottomrule
        \end{tabularx}
        \caption{Second preference pair ($j = 2$) of the two dialogue contexts in Figure~\ref{fig:context_ranking}. The dialogue context (student error type) and the corresponding utility $u(m, e)$ are shown above each pair. Posterior and utility values are listed in the order \textit{Focus}, \textit{Generic}, \textit{Probing}, \textit{Telling}; the score is their inner product.}
        \label{tab:example_pairs}
    \end{table}

    \paragraph{The Same Pattern in Every Error Type}
    Figure~\ref{fig:app_contexts} in Appendix~\ref{app:contexts} repeats Figure~\ref{fig:context_ranking} for every error type, drawing the context whose first pair has the widest normalized margin of its error type, with one row of blocks per family of the utility table. The pattern of Figure~\ref{fig:context_ranking} recurs in all seven contexts, and it follows the utility table. In the three \textit{Focus}-associated error types (first row of Figure~\ref{fig:app_contexts}), every candidate whose top label is \textit{Telling} or \textit{Probing} lies in the rejected half, at the very bottom of the ranking: 5 \textit{Telling} and 1 \textit{Probing} candidates for \textit{numerical calculation error}, 15 \textit{Probing} candidates for \textit{none of the above}, and 3 \textit{Telling} and 1 \textit{Probing} candidates for \textit{reached the correct solution but proceeded further}; the chosen halves of these three contexts consist of \textit{Focus} responses only, ordered by their \textit{Focus} posterior. In the three \textit{Telling}-associated error types (second row) the order is reversed: every candidate whose top label is \textit{Telling} stands at the top of the ranking and belongs to the chosen half (1, 12, and 11 candidates, with scores of 0.41, 1.35--1.50, and 0.57--1.05), and the rejected halves consist of \textit{Focus} responses whose scores differ only in the third decimal. For \textit{misunderstanding of the question} (third row), the only positive utility is that of \textit{Generic}, and the scores of the context stay below 0.07 with a nearly flat ranking after the first rank.

    \paragraph{Chosen--Rejected Probability Comparison}
    Figure~\ref{fig:posterior_shift} in Appendix~\ref{app:posterior_comparison} extends the comparison from the two contexts to all 24,298 pairs. For each error type and each move, it compares the mean classifier posterior of the chosen responses with that of the rejected responses. In all six error types with a \textit{Focus} or \textit{Telling} utility, the differences follow the utility table: the move with the largest utility is higher on the chosen side, and no move with zero utility is. For the three \textit{Focus}-associated error types, the mean \textit{Focus} posterior is 0.88--0.90 on the chosen side against 0.68--0.76 on the rejected side. For the three \textit{Telling}-associated error types, \textit{Telling} is higher on the chosen side (0.11--0.13 against 0.02), but the largest utility is not reflected in the largest posterior mass: \textit{Focus} remains the most probable move of the chosen responses (0.69--0.78). The utility table is thus reflected in the difference between the two sides, not in the dominant move of the chosen responses. \textit{Misunderstanding of the question} is the exception: its only positive utility belongs to \textit{Generic}, which the candidates almost never realize, so the two sides differ only marginally (0.004 against 0.002) and the scoring has nothing to prefer.

    Two properties of the data limit the \textit{Telling} side. First, the candidate pool is nearly the same for every error type and dominated by \textit{Focus}, the top label of 45,104 of the 48,839 candidates; \textit{Telling} is the top label of 1,498, and 307 of the 693 contexts hold no \textit{Telling}-labelled candidate at all (Table~\ref{tab:candidate_pool} in Appendix~\ref{app:posterior_comparison}). Second, the moves of the candidates are those assigned by the classifier of Section~\ref{subsec:selected_classifier}, which recovers about half of the \textit{Telling} utterances of its test set (F1 of 0.40), and the \textit{Telling} utilities themselves rest on one to four first tutor turns (Figure~\ref{fig:associations}). Both points are discussed in Chapter~\ref{chapter:Discussions}.

    \paragraph{Score Margins of the Constructed Pairs}
    Figure~\ref{fig:margin_by_pair} shows the score margins that the pairing rule of Section~\ref{subsec:preference_pair_construction} produces. Because pair \(j\) joins rank \(j\) with rank \(n-j+1\), the margin decreases with \(j\) in every context. The figure reports the mean normalized margin \(\Delta_{e,j}/\max_{m}\operatorname{PPMI}(m,e)\) over all contexts of each error type for the first 30 pairs of a context, the number of pairs that every context has. For the first pair, the mean lies between 0.37 (\textit{missing or incorrect factual knowledge}) and 0.71 (\textit{numerical calculation error} and \textit{none of the above}) in six of the seven error types; it is 0.03 for \textit{misunderstanding of the question}, where the three other utilities are zero and the scores of all candidates are bounded by the \textit{Generic} utility of 0.07 (third row of Figure~\ref{fig:app_contexts}). By the tenth pair the mean has fallen to 0.11--0.28, and by the thirtieth pair to 0.04 or less in every error type. The two error types with the largest \textit{Focus} utilities (\textit{numerical calculation error}, 0.90; \textit{none of the above}, 0.30) keep the widest margins at every \(j\); the other four error types with a \textit{Focus} or \textit{Telling} utility lie between 0.37 and 0.49 at the first pair and between 0.11 and 0.17 at the tenth.

    \begin{figure}[!t]
        \centering
        \includegraphics[width=\textwidth]{figures/fig_5_11_margin_by_pair.pdf}
        \caption[Score margins of the constructed pairs by pair index.]{Score margin of the constructed pairs by pair index \(j\) within a context (pair \(j\) joins rank \(j\) with rank \(n-j+1\)). Each point is the mean normalized margin \(\Delta_{e,j}/\max_{m}\operatorname{PPMI}(m,e)\) over all contexts of the error type. Only the first 30 pairs of every context are shown, so that every point rests on all contexts of its error type (87, 60, 49, 167, 97, 34, and 199 contexts in the order of the legend).}
        \label{fig:margin_by_pair}
    \end{figure}

    \paragraph{Generation Model of the Chosen and Rejected Responses}
    The six generation models contribute nearly equal shares of the candidate pool (15.7--17.0\,\% each) and, over all 24,298 pairs, nearly equal shares of the chosen side (12.1--19.5\,\%). Within an error type, however, the two sides differ in their model composition, and they differ in opposite directions in the two families (Figure~\ref{fig:model_composition}). For the three \textit{Telling}-associated error types, \path{gpt-5.6-terra} generates 27\,\% of the chosen but 5\,\% of the rejected responses, and \path{gemini-3.5-flash-lite} 6\,\% of the chosen but 26\,\% of the rejected responses; for the three \textit{Focus}-associated error types the picture is reversed, with \path{gpt-5.6-terra} at 5\,\% of the chosen and 26\,\% of the rejected responses. For \textit{misunderstanding of the question}, the shares of every model stay within 13--22\,\% on both sides. The pattern follows the moves that the models realize: \textit{Telling} is the top label of 12.5\,\% of the responses of \path{gpt-5.6-terra} but of at most 4\,\% of the responses of the other five models, so its responses rise to the chosen side wherever \textit{Telling} has a positive utility and fall to the rejected side wherever it has none. What this entails for preference optimization is discussed in Chapter~\ref{chapter:Discussions}.

    \begin{figure}[!t]
        \centering
        \includegraphics[width=\textwidth]{figures/fig_5_12_model_composition.pdf}
        \caption{Generation model of the chosen responses per error type. Each cell gives the share of the chosen responses of an error type that were generated by the model, with the model's share of all candidates of that error type (\textit{pool}) below; the colour shows the difference between the two in percentage points. Rows are labelled with the move that has the largest utility for the error type.}
        \label{fig:model_composition}
    \end{figure}

    \paragraph{The Preference Signal in the Constructed Dataset}
    Taken together, the figures answer the second part of Research Question~1. First, in each of the eight contexts shown (Figures~\ref{fig:context_ranking} and~\ref{fig:app_contexts}), the candidates that realize the move associated with the error type stand at the top of the ranking and become chosen, whereas the candidates that realize a move with zero utility stand at the bottom and become rejected. Second, the direction of this preference reverses between the \textit{Focus}-associated and the \textit{Telling}-associated error types: the same kind of response, a complete calculation, is rejected for \textit{numerical calculation error} and chosen for \textit{missing or incorrect factual knowledge} (Table~\ref{tab:example_pairs}). Third, the same holds on average over all 24,298 pairs: in every error type the chosen responses carry more posterior mass on the move with the largest utility than the rejected responses, and in the six error types with a \textit{Focus} or \textit{Telling} utility no move with zero utility carries more mass on the chosen side (Figure~\ref{fig:posterior_shift}). Fourth, the separation between the two halves is measurable: over all 24,298 pairs, the first pair of a context has a mean normalized margin between 0.37 and 0.71 in the six error types with a \textit{Focus} or \textit{Telling} utility (Figure~\ref{fig:margin_by_pair}). The constructed dataset therefore carries an error-conditioned preference signal for these six error types. For \textit{misunderstanding of the question}, whose only positive utility is that of \textit{Generic}, the scores are bounded by 0.07 and the margins stay close to zero (Figure~\ref{fig:margin_by_pair} and the third row of Figure~\ref{fig:app_contexts}); this case, and the fact that the signal is a difference in posterior mass while \textit{Focus} remains the most probable move of most responses on both sides (Table~\ref{tab:candidate_pool}), are taken up in Chapter~\ref{chapter:Discussions}.
